%% kapa: exam=1 round=27 year=2016 period=1 form=A subject=real_estate subject_ko=부동산학원론 seq=5 number=85
%% source: Q-Net 감정평가사 기출문제 (cst00309) artlSeq=5169402 file=extracted/1008/5169402/제27회 감정평가사 1차 문제/제27회 감정평가사 1교시문제(A형).pdf
%% conversion: kordoc 4.24.51 -> tools/gen1.py (auto); review=proofread-sonnet; flags=none
\begin{question}{85}
에스크로우(Escrow)제도에 관한 설명으로 옳지 않은 것은?
\begin{choices}
  \item 매수자는 권원상의 하자나 부담으로부터 발생하는 위험을 사전에 방지할 수 있다.
  \item 매수자뿐만 아니라 권원의 이전에 관계되는 매도자, 저당대출기관 등의 권익을 보호하는 역할을 한다.
  \item 권리보험제도와 병행하여 활성화하면 거래안전의 시너지 효과를 거둘 수 있다.
  \item 공인중개사법령상 개업공인중개사는 거래의 안전을 보장하기 위하여 필요하다고 인정하는 경우에는 거래계약의 이행이 완료될 때까지 계약금ㆍ중도금 또는 잔금을 개업공인중개사 명의로 금융기관에 예치하도록 거래당사자에게 권고할 수 있다.
  \item 에스크로우 회사는 매도자와 매수자의 협상과정에 참여하여 거래과정에서 발생하는 여러 가지 문제에 대하여 조언을 한다.
\end{choices}
\end{question}
\answer{5}
